\chapter{Complexity — Det Mod Math}
%%% generated by the blueprint pipeline from TCSlib.Complexity.ClassP.BinaryArith.DetModMath
\section{Overview}

Declarations of the Lean module \texttt{TCSlib.Complexity.ClassP.BinaryArith.DetModMath}.

\section{Declarations}

\begin{definition}[Matrix from row lists]
\lean{Complexity.BinArith.matOfRows}
\label{Complexity.BinArith.matOfRows}
\leanok
For natural numbers $p$ and $m$ and a list $R$ of lists of natural numbers, define an
$m\times m$ matrix over $\mathbb Z/p\mathbb Z$ by setting its $(i,j)$ entry to the
residue class of the $j$th entry of the $i$th row of $R$, using $0$ whenever either
entry is absent.
\end{definition}

\begin{definition}[Row elimination update]
\lean{Complexity.BinArith.elimRow}
\label{Complexity.BinArith.elimRow}
\leanok
For natural numbers $p$ and $\mathrm{inv}$ and lists $\mathrm{piv}$ and $\mathrm{row}$,
form a list by pairing the entries of the tails of $\mathrm{row}$ and $\mathrm{piv}$ up
to the length of the shorter tail. For each paired entry $(r,y)$, the output is the
remainder modulo $p$ of $r+p-(((\mathrm{head}(\mathrm{row})\cdot\mathrm{inv})\bmod
p)\cdot y\bmod p)$; the head of an empty row is taken to be $0$.
\end{definition}

\begin{theorem}[Indexing a List Tail]
\lean{Complexity.BinArith.getD_tail}
\label{Complexity.BinArith.getD_tail}
\leanok
\difficulty{4}
For any list $l$ of natural numbers and any natural number $j$, the element of the tail
of $l$ at index $j$, using $0$ as the default when that index is out of range, equals
the element of $l$ at index $j+1$ with the same default.
\end{theorem}

\begin{theorem}[The First Defaulted List Element]
\lean{Complexity.BinArith.getD_zero_eq_headD}
\label{Complexity.BinArith.getD_zero_eq_headD}
\leanok
\difficulty{4}
For any list $l$ of natural numbers, its element at index $0$, using $0$ as the default
for an empty list, equals its head, using $0$ when the list is empty.
\end{theorem}

\begin{theorem}[Indexing a Pointwise List Combination]
\lean{Complexity.BinArith.getD_zipWith}
\label{Complexity.BinArith.getD_zipWith}
\leanok
\difficulty{3}
Let $f$ be a function from pairs of natural numbers to natural numbers, and let
$l_1,l_2$ be lists of natural numbers. If $j$ is less than the length of each list, then
the element at index $j$ of the list formed by applying $f$ pointwise to $l_1$ and
$l_2$, with $0$ as the default out of range, equals $f$ applied to their respective
elements at index $j$, each read with default $0$.
\end{theorem}

\begin{theorem}[Subtraction in modular arithmetic]
\lean{Complexity.BinArith.natCast_sub_mod}
\label{Complexity.BinArith.natCast_sub_mod}
\leanok
\difficulty{3}
Let $p$ be a positive natural number, and let $r$ and $x$ be natural numbers. The
residue modulo $p$ of the remainder of $r+p-(x\bmod p)$ upon division by $p$ equals the
residue of $r$ minus the residue of $x$ in $\mathbb{Z}/p\mathbb{Z}$.
\end{theorem}

\begin{theorem}[Modular inverse by exponentiation]
\lean{Complexity.BinArith.natCast_pow_mod_eq_inv}
\label{Complexity.BinArith.natCast_pow_mod_eq_inv}
\leanok
\difficulty{4}
Let $p$ be prime and let $a$ be a natural number whose residue modulo $p$ is nonzero.
The residue modulo $p$ of the remainder of $a^{p-2}$ upon division by $p$ equals the
multiplicative inverse of the residue of $a$.
\end{theorem}

\begin{theorem}[Determinant under pivot elimination]
\lean{Complexity.BinArith.det_matOfRows_elim}
\label{Complexity.BinArith.det_matOfRows_elim}
\leanok
\uses{Complexity.BinArith.elimRow, Complexity.BinArith.getD_tail, Complexity.BinArith.getD_zero_eq_headD, Complexity.BinArith.getD_zipWith, Complexity.BinArith.matOfRows, Complexity.BinArith.natCast_sub_mod}
\difficulty{7}
Let $p$ be prime, let $v$ be a list of natural numbers and let $R$ be a list of rows,
with $v$ having length $|R|+1$ and every row in $R$ having length $|R|+1$. Suppose the
first entry of $v$, reduced modulo $p$, is nonzero, and let $u$ be a natural number
whose residue modulo $p$ is the inverse of that first entry. Then the determinant over
$\bbr/p\bbr$ of the matrix with rows $v$ followed by $R$ equals the first entry of $v$
modulo $p$ times the determinant of the matrix formed from the rows obtained by
eliminating the pivot using $u$.
\end{theorem}

\begin{theorem}[Determinant under row rotation]
\lean{Complexity.BinArith.det_matOfRows_rotate}
\label{Complexity.BinArith.det_matOfRows_rotate}
\leanok
\uses{Complexity.BinArith.matOfRows}
\difficulty{5}
For every natural number $p$, row $r_0$, and list of rows $R$, the determinant modulo
$p$ of the matrix with rows $R$ followed by $r_0$ equals $(-1)^{|R|}$ times the
determinant of the matrix with rows $r_0$ followed by $R$.
\end{theorem}

\begin{theorem}[Zero first column forces zero determinant]
\lean{Complexity.BinArith.det_matOfRows_eq_zero}
\label{Complexity.BinArith.det_matOfRows_eq_zero}
\leanok
\uses{Complexity.BinArith.getD_zero_eq_headD, Complexity.BinArith.matOfRows}
\difficulty{1}
Let $p,m$ be natural numbers and let $R$ be a list of natural-number rows. If, for every
$i<m+1$, the first entry of row $i$ (taken as zero when absent) is zero modulo $p$, then
the determinant modulo $p$ of the matrix of size $m+1$ formed from $R$, with missing
entries taken as zero, is zero.
\end{theorem}

\begin{theorem}[Matrix invariance under reduction]
\lean{Complexity.BinArith.matOfRows_map_mod}
\label{Complexity.BinArith.matOfRows_map_mod}
\leanok
\uses{Complexity.BinArith.matOfRows}
\difficulty{3}
For natural numbers $p,m$ and a list $R$ of lists of natural numbers, the $m\times m$
matrix over $\mathbb Z/p\mathbb Z$ obtained from $R$ is unchanged if every entry of $R$
is first reduced modulo $p$.
\end{theorem}

\begin{theorem}[Determinant of the empty matrix]
\lean{Complexity.BinArith.det_matOfRows_zero}
\label{Complexity.BinArith.det_matOfRows_zero}
\leanok
\uses{Complexity.BinArith.matOfRows}
\difficulty{1}
For every natural number $p$ and every list $R$ of lists of natural numbers, the
determinant of the $0\times 0$ matrix over $\mathbb{Z}/p\mathbb{Z}$ whose entries are
read from $R$ (using zero when an entry is missing) is $1$.
\end{theorem}
